\documentclass{ximera}

\title{Project 3: Jada Ramey}
\author{Jada Ramey}

\begin{document}
\begin{abstract}
    Project 3 on kaleidoscope completed by Jada Ramey.
\end{abstract}

\maketitle

Note: The original template could not be accessed due to editor limits, so this document was created independently in Overleaf while following the project requirements.

\section*{Introduction}
This project explores reflections and rotations using matrices to simulate a kaleidoscope. Linear transformations are used to model symmetry and repeated patterns.

\section*{(1) Reflection across the x-axis}

\[
\begin{bmatrix}
1 & 0 \\
0 & -1
\end{bmatrix}
\]

This matrix reflects points across the x-axis by keeping the x-coordinate the same and negating the y-coordinate.

\section*{(2) Rotation by 60 degrees}

\[
\begin{bmatrix}
\frac{1}{2} & -\frac{\sqrt{3}}{2} \\
\frac{\sqrt{3}}{2} & \frac{1}{2}
\end{bmatrix}
\]

This matrix rotates points counterclockwise by $60^\circ$.

\section*{(3) Reflection across a 60 degree line}

\[
\begin{bmatrix}
-\frac{1}{2} & \frac{\sqrt{3}}{2} \\
\frac{\sqrt{3}}{2} & \frac{1}{2}
\end{bmatrix}
\]

This matrix was constructed using the reflection formula for a line at a 60 degree angle.

\section*{(4a) Composition}

Let $R$ be the rotation matrix and $F$ be the reflection matrix.

\[
R =
\begin{bmatrix}
\frac{1}{2} & -\frac{\sqrt{3}}{2} \\
\frac{\sqrt{3}}{2} & \frac{1}{2}
\end{bmatrix},
\quad
F =
\begin{bmatrix}
1 & 0 \\
0 & -1
\end{bmatrix}
\]

\[
RF \neq FR
\]

Matrix multiplication is not commutative, so the order of transformations matters. Applying reflection first versus rotation first produces different results.

\section*{(4b) Powers of R}

\[
R^6 =
\begin{bmatrix}
1 & 0 \\
0 & 1
\end{bmatrix}
\]

Applying the rotation six times results in a full $360^\circ$ rotation, returning the image to its original position.

\section*{(4c) Determinants}

\[
\det(R) = 1, \quad \det(F) = -1
\]

A determinant of 1 preserves orientation, while a determinant of -1 reverses orientation. This explains why reflections flip images while rotations do not.

\end{document}